% Generated draft: requires mathematical and source review.
\documentclass[11pt]{article}
\usepackage[margin=25mm]{geometry}
\usepackage{fontspec}
\setmainfont{Latin Modern Roman}
\usepackage{amsmath,amssymb,mathtools,unicode-math}
\setmathfont{Latin Modern Math}
\usepackage{enumitem,needspace,xurl,hyperref}
\setlength{\parindent}{0pt}
\setlength{\parskip}{5pt}
\setlength{\emergencystretch}{3em}
\newcommand{\N}{\mathbb{N}}
\newcommand{\Z}{\mathbb{Z}}
\newcommand{\Q}{\mathbb{Q}}
\newcommand{\R}{\mathbb{R}}
\newcommand{\C}{\mathbb{C}}
\newcommand{\F}{\mathbb{F}}
\newcommand{\A}{\mathbb{A}}
\newcommand{\PP}{\mathbb{P}}
\newcommand{\E}{\mathbb{E}}
\newcommand{\eps}{\varepsilon}
\newcommand{\dd}{\,\mathrm{d}}
\begin{document}
\section*{UnsolvedMath ID: 30001458}
\textbf{Source problem code:} catalogue entry
\subsection{Definitions and mathematical statement}
Let \(K_n\) denote the polygonal approximations in the Koch-snowflake construction, and let \(K\) denote their limiting snowflake domain. The supplied record does not specify the indexing, normalization, or precise surface construction associated with these approximations.

Under the usual translation-surface interpretation, let \((X_n,\omega_n)\) be a chosen unfolding of the rational polygon \(K_n\). A translation surface has coordinate changes given by translations away from its singularities. Its Veech group is
\[
\Gamma_n=\{D f: f\in\operatorname{Aff}^{+}(X_n,\omega_n)\}\subseteq\operatorname{SL}_2(\R),
\]
where \(\operatorname{Aff}^{+}\) denotes the orientation-preserving affine self-homeomorphisms in translation coordinates. This interpretation supplies terminology; the record does not identify the particular unfoldings or any compatible maps between them.

An inverse limit of these groups requires specified homomorphisms \(p_n:\Gamma_{n+1}\to\Gamma_n\). Given such maps, its underlying group would be
\[
\varprojlim(\Gamma_n,p_n)
=\left\{(g_n)_{n\geq 0}\in\prod_{n\geq 0}\Gamma_n:
 p_n(g_{n+1})=g_n\text{ for every }n\geq 0\right\}.
\]
No bonding homomorphisms are supplied, so this expression does not yet define a particular group. A topology would also need to be specified if a topological or dynamical role is intended.

\textbf{Source question, with unresolved qualifications retained.} Does the inverse limit of the Veech groups of the prefractal Koch-snowflake surfaces define a dynamically relevant object yielding a well-defined geodesic flow on a fractal flat surface dynamically equivalent to the Koch-snowflake billiard flow?

For this question to become a precise existence statement, one must specify the inverse system, the limiting surface and its flat or metric structure, its geodesic phase space and treatment of singularities, and a definition and domain of the billiard flow in \(K\). The ordinary polygonal reflection rule does not by itself specify reflection at points of the limiting boundary without a tangent, or continuation through singular collisions. One must also specify the sense of dynamical equivalence, including whether it preserves time and whether it holds everywhere or only on a specified invariant subset. The supplied record makes none of these choices, and no particular choice is asserted here. Thus the retained target is a conjectural construction program rather than a fully specified theorem or conjecture.
\subsection{Short English statement}
The problem asks whether a compatible limit of the Veech groups of Koch-snowflake approximations can produce a fractal flat surface whose geodesic dynamics correspond to billiard dynamics in the limiting snowflake. The connecting maps, limiting geometry, boundary dynamics, and meaning of equivalence remain unspecified. The supplied record labels the problem open; this draft does not independently verify that status.
\subsection{Sources}
\begin{itemize}
\item Supplied local UnsolvedMath registry record, permanent ID \(30001458\), source code \texttt{catalogue entry}, title \emph{Geodesic Flows on Koch Snowflake Fractal Surfaces}, in the Oberwolfach Reports Open Problems collection. The record gives proposed year 2010 and names no individual proposer. It labels the statement \texttt{exact}; its supplied verification says only that the internally readable wording was retained byte-for-byte.
\item Supplied report attribution: \emph{Arbeitsgemeinschaft: Mathematical Billards} (2010), \emph{Oberwolfach Reports}, pp.~955--1015, DOI \texttt{10.4171/owr/2010/17}. Supplied source URL and literature-source link, labeled ``original OWR report'': \url{https://doi.org/10.4171/owr/2010/17}. The background describes Oberwolfach Reports as documenting mathematical workshops and their open-problem discussions. No individual problem page or section is supplied.
\item Literature assessment dated exactly \texttt{2026-08-21}: the supplied curation labels the problem open and describes the inverse-limit Veech-group and fractal-flat-surface construction as a conjectural program. It reports that no primary source located by its cutoff gives a general construction with the claimed geodesic and billiard-flow equivalence. This is supplied dated curation, not an independently verified literature assessment; the record itself requests independent verification of time-sensitive status and cited claims.
\end{itemize}
\end{document}
